
\documentclass[11pt]{article}

\usepackage{amsmath,amssymb,amsthm,mathtools}
\usepackage[T1]{fontenc}
\usepackage{lmodern}
\usepackage[margin=1in]{geometry}
\usepackage{enumitem}
\usepackage{booktabs}
\usepackage{array}
\usepackage{longtable}
\usepackage{microtype}
\usepackage{xcolor}
\usepackage{hyperref}
\usepackage{xurl}

\hypersetup{colorlinks=true,linkcolor=blue,citecolor=blue,urlcolor=blue}
\urlstyle{same}

\newtheorem{theorem}{Theorem}[section]
\newtheorem{frontier}[theorem]{Frontier theorem}
\newtheorem{background}[theorem]{Background input}
\newtheorem{proposition}[theorem]{Proposition}
\newtheorem{lemma}[theorem]{Lemma}
\newtheorem{corollary}[theorem]{Corollary}
\newtheorem{remark}[theorem]{Remark}
\newtheorem{definition}[theorem]{Definition}
\newtheorem{heuristic}[theorem]{Heuristic}
\newtheorem{question}[theorem]{Open problem}

\newcommand{\Z}{\mathbb Z}
\newcommand{\sig}[2]{\sigma_{#1#2}}
\newcommand{\Lw}[1]{L_{#1}}
\newcommand{\Rw}[1]{R_{#1}}
\newcommand{\CanL}{[2,2,1]}
\newcommand{\CanR}{[1,4,4]}
\newcommand{\CanSeed}{\CanL/\CanR}
\newcommand{\CL}{\mathcal C_L}
\newcommand{\CR}{\mathcal C_R}
\newcommand{\Omix}{\mathcal O_{\mathrm{mix}}}
\newcommand{\Obr}{\mathcal O_{\mathrm{br}}}
\newcommand{\RhoR}{\rho_{R1}}
\newcommand{\SelCorr}{\mathrm{Sel}_{\mathrm{corr}}}
\newcommand{\SelStar}{\mathrm{Sel}^{\star}}
\title{Odd-$m$, $d=5$ as a one-cornered odometer:\\
consolidated theorem-order manuscript}
\author{}
\date{March 22, 2026}

\begin{document}
\maketitle

\begin{abstract}
This text consolidates the current odd-$m$, $d=5$ proof chain in the theorem
order most useful for later line-by-line cleanup and formalization.
The proof is organized around the spine
\[
\begin{aligned}
&\text{front-end one-corner reduction}
\to
\text{exact odometer spine}\\
&\to
\text{selector surgery}
\to
\text{generic upgrade and globalization}.
\end{aligned}
\]
Inside the present bundle, the front-end slice
\[
T0,T1,T2,T3,T4
\]
and the selector slice
\[
G1,G2
\]
are closed directly.
The remaining manuscript-order inputs are no longer left as anonymous imports:
the post-entry one-corner odometer theorem is tied to the compact one-pass
odd-$m$ globalization package, and the graph/globalization endgame is tied
explicitly to the accepted theorem-side anchor package, the corrected-selector
baseline theorem, the color-$4$ Sel$^{\star}$ Hamilton theorem, and the
odd-$m$ globalization theorem.
What still remains explicit at the honest-boundary level is one graph-side
transport compatibility input: the cyclic/color-relative hypothesis needed to
apply the abstract \(G2\) transport theorem to the concrete two-swap package
constructed in \(G1\).
With that caveat stated honestly, the final theorem reduces to a short corrected
generic upgrade step.
The resulting manuscript is not a first-principles rebuild of every historical
artifact, but it is a closed theorem-order document: every load-bearing theorem
used in the final proof is now named, stated, and placed in the dependency
chain.
\end{abstract}

\tableofcontents

\section{The global proof spine}

The intended theorem-level structure of the odd-$m$, $d=5$ paper is
\[
\begin{aligned}
&\text{front-end one-corner reduction}
\to
\text{exact odometer spine}\\
&\to
\text{selector surgery}
\to
\text{generic upgrade and globalization}.
\end{aligned}
\]
At a slightly finer level the manuscript is organized as follows.
\begin{enumerate}[label=\textup{(\Roman*)},leftmargin=3.0em]
\item \textbf{Front-end one-corner reduction.}
      Reduce every unresolved front-end instance to the canonical bridge seed
      \(\CanSeed\) by a finite chart plus three shape-side row theorems.
\item \textbf{Canonical entry and exact odometer spine.}
      Show that the canonical seed has a unique live unresolved continuation,
      namely the alternate-$2$ entry
      \[
      E=(m-1,1,u_{\mathrm{source}},2),
      \]
      and then identify the resulting dynamics with the exact one-corner
      odometer corridor.
\item \textbf{Selector surgery.}
      Insert the already closed color-$4$ Hamilton branch into a colorwise
      selector package by a selector-compatibility theorem.  The remaining
      colors follow once that explicit package is tied to the cyclic transport
      hypothesis.
\item \textbf{Final stitching and globalization.}
      Use the colorwise graph package together with the explicit odometer block
      to produce the final odd-$m$, $d=5$ Hamilton decomposition.
\end{enumerate}

\begin{theorem}[actual-needed tiny-repair core]
\label{thm:T0-needed}
For every odd modulus $m>2$, the front-end uses of tiny repair that actually
occur below are already closed internally.
More precisely:
\begin{enumerate}[label=\textup{(\arabic*)},leftmargin=2.8em]
\item at
      \[
      \sig21=[4,1,4]/[1,2,2],
      \]
      every one-label one-gate or one-bit repair branch supported on
      \(L2,L3,R2,\) or \(R3\), with cocycle-defect choice in
      \(\{\mathrm{none},\mathrm{left},\mathrm{right},\mathrm{both}\}\),
      is non-Latin;
\item at the canonical bridge seed \(\CanSeed\), every one-label one-gate or
      one-bit repair branch supported on \(L3,R2,\) or \(R3\), with the same
      four cocycle-defect choices, is non-Latin.
\end{enumerate}
\end{theorem}

\begin{remark}[scope convention for ``after tiny repair'']
Throughout the front-end section, the phrase ``after tiny repair'' refers only
to the concrete removals supplied by Theorem~\ref{thm:T0-needed} at the relevant
seed.  It does not silently quantify over richer multi-class repair families.
A stronger standalone one-label repair theorem for the full artifact seed list
may still be desirable later, but it is not used below.
\end{remark}

\begin{theorem}[main odd-$m$, $d=5$ Hamilton decomposition theorem, conditional on transport compatibility]
\label{thm:conditional-full}
Assume that the explicit package of Theorem~\ref{fr:G1} is obtained from one
color-relative local rule family in the sense of
Theorem~\ref{thm:G2-closed}.
Then for every odd modulus $m\ge 5$, the directed torus
\[
D_5(m)=\operatorname{Cay}((\Z/m\Z)^5,\{e_0,e_1,e_2,e_3,e_4\})
\]
admits a decomposition into five arc-disjoint directed Hamilton cycles.
\end{theorem}

The point of the present manuscript is now slightly different from the earlier
``honest status'' draft.
The front-end slice $T0$--$T4$ and the selector slice $G1,G2$ are already
closed in this bundle, and the remaining theorem packets used by the final
proof are recorded below as explicit named upstream theorems.
This makes the dependency chain linear enough for later polishing and for a
Lean-oriented theorem ledger.

\section{Three languages and a methodological rule}

\begin{remark}[support language, grid language, odometer language]
Three languages appear repeatedly and should not be confused.
\begin{enumerate}[label=\textup{(\arabic*)},leftmargin=2.8em]
\item \textbf{Support language.}
      This records only which local families are present.
      It is a pruning language.
\item \textbf{Grid language.}
      This records actual endpoint words, equivalently the grid cell
      \(\sig{i}{j}\).
      It is the main language of shape-side row theorems.
\item \textbf{Odometer language.}
      This records the explicit post-entry coordinates
      \((q,w,u,\Theta)\).
      It is the main language after the canonical entry theorem.
\end{enumerate}
The key rule is:
\emph{support is a pruning invariant, not a classification invariant.}
So support can help isolate a theorem target, but it cannot by itself separate
all rows that the grid language can separate.
\end{remark}

\section{Finite front-end chart}

\subsection{\texorpdfstring{The \(3\times 3\)}{The 3x3} bridge-position grid}

Define
\[
L_1=[1,4,4],\qquad L_2=[4,1,4],\qquad L_3=[4,4,1],
\]
and
\[
R_1=[1,2,2],\qquad R_2=[2,1,2],\qquad R_3=[2,2,1].
\]
The noncanonical rows are arranged in the grid
\[
\sig{i}{j}=L_i/R_j,\qquad 1\le i,j\le 3.
\]
The interpretation is:
\begin{itemize}[leftmargin=2.4em]
\item \(i\) records the position of the bridge symbol on the left endpoint;
\item \(j\) records the position of the bridge symbol on the right endpoint.
\end{itemize}

\begin{proposition}[finite seed certificate]
\label{prop:finite-chart}
The filtered front-end seed set contains exactly eleven raw rows.
Under the accepted symmetries, these reduce to five classes
\[
C0,\ C1,\ C2,\ C3,\ C4.
\]
Exactly one class is an oriented bridge class, namely \(C0\), and it contains
the canonical bridge seed
\[
\CanSeed=\CanL/\CanR.
\]
The centered-bridge classes are exactly
\[
C1\cup C3\cup C4,
\]
and the residual nonbridge class outside that layer is exactly
\[
C2=\sig33=[4,4,1]/[2,2,1].
\]
\end{proposition}

\begin{proposition}[bridge normalization]
\label{prop:bridge-normalization}
Every oriented bridge row is equivalent under the accepted symmetries to the
canonical seed
\[
\CanSeed=[2,2,1]/[1,4,4].
\]
\end{proposition}

\begin{definition}[changed-family support]
For each surviving row \(\sigma\), let
\[
\Sigma(\sigma)\subseteq\{L1,L2,L3,R1,R2,R3\}
\]
denote the changed-family support.
On the noncanonical \(3\times 3\) grid the support factorization is
\[
\Sigma(\sig{i}{j})
=
\{L2,L3,R1,R3\}
\cup (\{L1\}\ \text{if } i=1)
\cup (\{R2\}\ \text{if } j\le 2).
\]
\end{definition}

\begin{remark}[the exceptional row]
The row
\[
\sig11=[1,4,4]/[1,2,2]
\]
is exceptional for a finite reason:
it is the unique bridge-first row in which the left bridge is still leading.
In the bridge-position grid it is the double-leading corner.
\end{remark}

\subsection{Prototype reduction on the shape side}

\begin{proposition}[two-prototype reduction]
\label{prop:two-prototypes}
Under the accepted symmetries, the centered-bridge layer decomposes into two
orbits:
\[
\Omix=
\{\sig21,\ \sig23,\ \sig12,\ \sig32\},
\qquad
\{\sig22\}.
\]
Thus the whole centered-bridge burden reduces to one mixed prototype
\(\sig21=[4,1,4]/[1,2,2]\) and one double-centered core \(\sig22\).
\end{proposition}

\begin{corollary}[shape-side compression]
\label{cor:shape-side-compression}
Modulo Theorem~\ref{thm:T0-needed} and bridge normalization, the front-end shape side
reduces to exactly three row theorem objects:
\[
T1:\ \sig21\rightsquigarrow \sig31,\qquad
T2:\ \sig22\text{ closes},\qquad
T3:\ \sig33\text{ closes}.
\]
\end{corollary}

\section{Closed front-end row theorems}
\label{sec:frontier}

\begin{theorem}[T1: first-column companion theorem]
\label{fr:T1}
After the \(\sig21\)-repair branches from Theorem~\ref{thm:T0-needed} are
removed, every live unresolved immediate continuation of
\[
\sig21=[4,1,4]/[1,2,2]
\]
is represented by the bridge companion
\[
\sig31=[4,4,1]/[1,2,2].
\]
\end{theorem}

\begin{theorem}[T2: double-centered closure]
\label{fr:T2}\label{thm:T2-closed}
For every odd modulus \(m>2\), no row obtained from
\[
\sig22=[4,1,4]/[2,1,2]
\]
by a one-gate or one-bit one-label repair, with cocycle-defect choice in
\(\{\mathrm{none},\mathrm{left},\mathrm{right},\mathrm{both}\}\), is Latin.
Equivalently, after the artifact-surfaced tiny-repair scope is exhausted, the
center cell \(\sig22\) does not remain as a live unresolved front-end row.
\end{theorem}

\begin{theorem}[T3: repeated-end corner closure]
\label{fr:T3}\label{thm:T3-closed}
For every odd modulus \(m>2\), no row obtained from
\[
\sig33=[4,4,1]/[2,2,1]
\]
by a one-gate or one-bit one-label repair, with cocycle-defect choice in
\(\{\mathrm{none},\mathrm{left},\mathrm{right},\mathrm{both}\}\), is Latin.
Equivalently, after the artifact-surfaced tiny-repair scope is exhausted, the
repeated-end corner \(\sig33\) does not remain as a live unresolved front-end row.
\end{theorem}

\begin{remark}[proof packets for \(T2\) and \(T3\)]
The odd-\(m\ge 5\) obstruction proofs are extracted in the companion notes
\path{d5_228_T2_centered_core_obstruction_note.tex} and
\path{d5_225_T3_four_collision_families_note.tex}.
The remaining one-label \(m=3\) cases are exhausted by the finite scans recorded
in \path{d5_229_T2_T3_m3_one_label_scan.py} and its saved summaries.
\end{remark}

\begin{proposition}[canonical support certificate]
\label{prop:canonical-support}
At the canonical bridge seed
\[
\CanSeed=\CanL/\CanR,
\]
the support is exactly
\[
\Sigma(\CanSeed)=\{L3,R1,R2,R3\}.
\]
\end{proposition}

\begin{theorem}[T4: canonical entry theorem]
\label{fr:T4}
At the canonical bridge seed \(\CanSeed\), after the three dead repair source
families from Theorem~\ref{thm:T0-needed}---namely \(L3\), \(R2\), and
\(R3\)---are removed, the unique live immediate source family is \(R1\).
Its canonical two-step normalization in the shifted section coordinates
\[
(Q,W,U,\Theta)=(q-1,\ w+1,\ u-1,\ 2)
\]
is the alternate-$2$ entry
\[
E=(m-1,1,u_{\mathrm{source}},2).
\]
\end{theorem}

\begin{theorem}[front-end one-corner reduction]
\label{thm:conditional-front-end}
Every unresolved odd-$m$, $d=5$ front-end instance reduces to the exact
one-corner channel with canonical seed \CanSeed and entry state
\[
E=(m-1,1,u_{\mathrm{source}},2).
\]
\end{theorem}

\begin{proof}
By Proposition~\ref{prop:finite-chart}, every unresolved front-end instance is
represented by one of the finitely many rows in the chart.
Proposition~\ref{prop:bridge-normalization} moves every oriented bridge row to
the canonical seed \CanSeed.  Corollary~\ref{cor:shape-side-compression}
reduces the noncanonical shape side to the theorem objects \(T1,T2,T3\),
and Theorems~\ref{fr:T1}--\ref{fr:T4} close those row objects and identify
the unique live bridge continuation with the entry state \(E\).  Hence every
unresolved front-end instance reduces to the exact one-corner channel with
canonical seed \CanSeed and entry state \(E\).
\end{proof}

\section{\texorpdfstring{The closure package for \(T1\)}{The closure package for T1}}

What used to be the live front-end proof search is now a closed package.
We keep the strip analysis here because it explains why Theorem~\ref{fr:T1} is
short.

\subsection{Unique surviving branch}

\begin{proposition}[source support for \(T1\)]
\label{prop:T1-source-support}
At the mixed prototype
\[
\sig21=[4,1,4]/[1,2,2]
\]
one has
\[
\Sigma(\sig21)=\{L2,L3,R1,R2,R3\}.
\]
By Theorem~\ref{thm:T0-needed}, every immediate branch except the
\(R1\)-branch is removed.
\end{proposition}

So \(T1\) is not a many-branch problem.
It is a unique surviving-branch problem.

\subsection{Two strips and two local raw targets}

Define the centered-left strip and centered-right strip by
\[
\CL=\{\sig21,\sig22,\sig23\},
\qquad
\CR=\{\sig12,\sig22,\sig32\}.
\]
With Theorems~\ref{thm:T2-closed} and \ref{thm:T3-closed} recorded as
closed obstruction cells, \(T1\) reduces to excluding the dangerous targets
on these two strips.

\begin{proposition}[left-strip case split]
\label{prop:left-strip-case-split}
Relative to the source right word \([1,2,2]\), the centered-left strip has
exactly three cases:
\[
\sig21 \leftrightarrow \Delta_R=\varnothing,\qquad
\sig22 \leftrightarrow \Delta_R=\{1,2\},\qquad
\sig23 \leftrightarrow \Delta_R=\{1,3\}.
\]
In particular,
\[
\sig23=[4,1,4]/[2,2,1]
\]
is the unique centered-left case with disconnected change-set \(\{1,3\}\), the
unique centered-left case with terminal right-slot change, and the unique
centered-left case with \(R2\) absent.
\end{proposition}

\begin{proposition}[current \(T1\) reduction after strip closure]
\label{prop:T1-current-reduction}
Assume Theorems~\ref{thm:T2-closed} and \ref{thm:T3-closed}.
Then \(T1\) has no remaining live raw local target.
The left-strip target that the surviving immediate \(R1\)-image out of
\(\sig21\) does not realize \(\sig23\) is closed by
Theorem~\ref{thm:sigma23-closed}, and the normalized centered-right target that
it does not realize \(\sig32\) is closed by
Theorem~\ref{thm:sigma32-closed} below.
\end{proposition}

\subsection{The \texorpdfstring{$\sig23$}{sigma23} side is now closed}

The bookkeeping reductions from the earlier notes remain useful, but the live
frontier theorem on this side is no longer open.  Once the surfaced bundle
builder is admitted, the surviving immediate \(R1\)-family out of \(\sig21\)
can be read directly and its second-step comparison is explicit.

\begin{proposition}[equivalent forms of \(\sig23\) exclusion]
\label{prop:sigma23-equivalences}
For the surviving immediate \(R1\)-image \(\RhoR\) out of \(\sig21\), each of the
following suffices to exclude \(\sig23\).
\begin{enumerate}[label=\textup{(\arabic*)},leftmargin=2.8em]
\item \(\Delta_R(\RhoR)\subseteq\{1,2\}\);
\item the terminal right slot is unchanged;
\item the changed right slots form a connected set;
\item the immediate update is an elementary rewrite rooted at the acted edge
      \(\{1,2\}\);
\item the immediate update is supported on the closed star of the acted edge.
\end{enumerate}
\end{proposition}

\begin{theorem}[surfaced-bundle closure of the \(\sig23\) side]
\label{thm:sigma23-closed}
For every modulus \(m>2\), the actual second step of the surviving immediate
\(R1\)-family out of
\[
\sig21=[4,1,4]/[1,2,2]
\]
agrees with the \(\sig22=[4,1,4]/[2,1,2]\) prototype and never with the
\(\sig23=[4,1,4]/[2,2,1]\) prototype.
Equivalently, for the tiny-repair-normalized surviving immediate \(R1\)-image
\(\RhoR\) one has
\[
\Delta_R(\RhoR)=\{1,2\}.
\]
In particular the left-strip burden inside \(T1\) is closed.
\end{theorem}

\begin{proof}
Fix a modulus \(m>2\), color \(0\), and the standard normalization
\(w_0=s_0=0\).
Write \(\partial_i\) for the coordinate step that adds \(+1\) in direction
\(i\) modulo \(m\).
In the surfaced builder, the common right \(R1\)-source family for
\(\sig21,\sig22,\sig23\) is exactly
\[
S_R=\{x:\ \mathrm{layer}(x)=1,\ q(x)=m-1,\ w(x)=0,\ s(x)\neq 0\}.
\]
Thus all three rows are compared on the same source family.

For \(\sig21\), the right word is \([1,2,2]\), so the first two right-side
steps on \(S_R\) are
\[
x\mapsto \partial_1(x)\mapsto \partial_2\partial_1(x).
\]
For \(\sig22\), the right word is \([2,1,2]\), so the first two right-side
steps are
\[
x\mapsto \partial_2(x)\mapsto \partial_1\partial_2(x).
\]
For \(\sig23\), the right word is \([2,2,1]\), so the first two right-side
steps are
\[
x\mapsto \partial_2(x)\mapsto \partial_2^2(x).
\]

The surfaced common module defines each \(\partial_i\) as coordinatewise
\(+1\) in a single axis modulo \(m\).  Therefore the coordinate steps
commute:
\[
\partial_2\partial_1=\partial_1\partial_2.
\]
Thus the actual second step of \(\sig21\) on the surviving \(R1\)-family agrees
exactly with the \(\sig22\) prototype.

On the other hand, every \(x\in S_R\) satisfies \((q(x),w(x))=(m-1,0)\).
Therefore
\[
\partial_2\partial_1(x):(q,w)=(0,1),\qquad
\partial_2^2(x):(q,w)=(m-1,2).
\]
These are distinct for every \(m>2\).  Hence the actual second step of
\(\sig21\) does not realize \(\sig23\).
The description of \(\Delta_R\) now follows from
Proposition~\ref{prop:left-strip-case-split}.
\end{proof}

\subsection{The \texorpdfstring{$\sig32$}{sigma32} side is also closed}

The remaining centered-right target is the repeated-end option
\[
\sig32=[4,4,1]/[2,1,2].
\]
Here the common comparison family is on the left rather than on the right.
The first two left-side steps are again explicit in the surfaced builder.

\begin{theorem}[surfaced-bundle closure of the \(\sig32\) side]
\label{thm:sigma32-closed}
For every modulus \(m>2\), the normalized surviving immediate \(R1\)-image out
of
\[
\sig21=[4,1,4]/[1,2,2]
\]
does not realize the dangerous centered-right target
\[
\sig32=[4,4,1]/[2,1,2].
\]
More precisely, on the common left family
\[
S_L=\{x:\ \mathrm{layer}(x)=1,\ q(x)=m-1,\ w(x)=m-1,\ s(x)\neq 0\},
\]
the actual two-step left continuation agrees with the centered core
\(\sig22=[4,1,4]/[2,1,2]\) and differs from the repeated-end target
\(\sig32\).
\end{theorem}

\begin{proof}
Fix a modulus \(m>2\), color \(0\), and the standard normalization
\(w_0=s_0=0\).
On the common left family \(S_L\), the four relevant left words are
\[
\sig21:\ [4,1,4],\qquad
\sig12:\ [1,4,4],\qquad
\sig22:\ [4,1,4],\qquad
\sig32:\ [4,4,1].
\]
Therefore the first two left-side steps on \(S_L\) are
\[
\sig21:\ x\mapsto \partial_4(x)\mapsto \partial_1\partial_4(x),
\]
\[
\sig12:\ x\mapsto \partial_1(x)\mapsto \partial_4\partial_1(x),
\]
\[
\sig22:\ x\mapsto \partial_4(x)\mapsto \partial_1\partial_4(x),
\]
\[
\sig32:\ x\mapsto \partial_4(x)\mapsto \partial_4^2(x).
\]
Hence the actual left-side continuation of \(\sig21\) agrees with the
\(\sig22\) prototype.

To distinguish it from \(\sig32\), note that every \(x\in S_L\) satisfies
\(q(x)=m-1\).  Therefore
\[
\partial_1\partial_4(x): q=0,
\qquad
\partial_4^2(x): q=m-1.
\]
These are distinct for every \(m>2\).  Hence the normalized surviving
immediate \(R1\)-image out of \(\sig21\) does not realize \(\sig32\).
\end{proof}

\begin{remark}[what remains after strip closure]
The \(\sig23\)- and \(\sig32\)-sides are no longer frontier theorems.
Their remaining work is only packaging: the surfaced-bundle proofs above should
be translated into the final preferred notation of the main text.
Under Proposition~\ref{prop:T1-current-reduction}, there is now no remaining
live raw local burden inside \(T1\).
\end{remark}

\begin{proof}[Proof of Theorem~\ref{fr:T1}]
By Proposition~\ref{prop:T1-source-support} and Theorem~\ref{thm:T0-needed},
only the immediate \(R1\)-branch out of \(\sig21\) survives the actual front-end
repair cleanup.  Proposition~\ref{prop:T1-current-reduction}, together with
Theorems~\ref{thm:sigma23-closed} and \ref{thm:sigma32-closed}, shows that this
surviving branch does not realize either dangerous strip target.
Theorems~\ref{thm:T2-closed} and \ref{thm:T3-closed} remove the centered core
and repeated-end obstruction cells themselves.  Hence every live unresolved
immediate continuation of \(\sig21\) is represented by the bridge companion
\(\sig31\), as claimed.
\end{proof}

\subsection{\texorpdfstring{Why the present \(T1\) target is now exhausted}{Why the present T1 target is now exhausted}}

The moral is that \(T1\) is no longer a package jungle.
After the strip-closure theorems, it became a unique surviving-branch theorem
with no remaining live raw local comparison.  The only auxiliary inputs are the
already-closed obstruction theorems \(T2\) and \(T3\).

\section{Canonical entry and the exact odometer spine}

\subsection{The reduced entry problem}

\begin{proposition}[source support for \(T4\)]
\label{prop:T4-source-support}
At the canonical bridge seed
\[
\CanSeed=\CanL/\CanR,
\]
the only support families present are \(L3,R1,R2,R3\).
By Theorem~\ref{thm:T0-needed}, the branches \(L3\), \(R2\), and \(R3\) are
dead, so the only potentially live immediate continuation is the \(R1\)-branch.
\end{proposition}

Thus the surviving-branch toolkit is now fully closed on the front end, and
its output language is the odometer language rather than the grid language.

\subsection{The explicit odometer block}

With Theorem~\ref{fr:T4} proved, the rest of the theorem-side front end is explicit.

\begin{definition}[candidate one-corner orbit]
Fix odd \(m\) and the entry state
\[
E=(m-1,1,u_{\mathrm{source}},2).
\]
Let
\[
\rho:=u_{\mathrm{source}}+1 \pmod m.
\]
Define the candidate orbit \(\widehat x_n=(q_n,w_n,u_n,\Theta_n)\) by
\begin{align*}
\widehat F(q,w,u,0) &= (q+1,w,u,1),\\
\widehat F(q,w,u,1) &= (q,w,u+1,2),\\
\widehat F(q,w,u,2) &= (q,w+\mathbf 1_{q=m-1},u,3),\\
\widehat F(q,w,u,\Theta) &= (q,w,u,\Theta+1)\qquad(3\le \Theta\le m-1).
\end{align*}
On the section \(\Theta=2\), the induced return map is
\[
(q,w)\mapsto(q+1,w+\mathbf 1_{q=m-1}).
\]
\end{definition}

\begin{definition}[odometer coordinate]
Define
\[
\eta(q,w)=q+m(w-1)-(m-1).
\]
Then \(\eta\) increases by exactly \(1\) under each return to the section
\(\Theta=2\).
\end{definition}

\begin{lemma}[phase-$1$ residue invariant]
Along the candidate orbit, every phase-$1$ state satisfies
\[
q\equiv u-\rho+1 \pmod m.
\]
\end{lemma}

\begin{theorem}[post-entry exact odometer spine]
\label{thm:odometer-spine}
The actual active front-end branch agrees with the candidate one-corner
orbit up to first exit.
The trigger family is
\[
H_{L1}=\{(q,w,u,\Theta)=(m-1,m-1,u,2):u\neq 2\},
\]
and the universal first exits are
\[
T_{\mathrm{reg}}=(m-1,m-2,1),
\qquad
T_{\mathrm{exc}}=(m-2,m-1,1).
\]
Moreover, no patched current class appears before exit.
\end{theorem}

\begin{proof}
Theorem~\ref{fr:T4} identifies the unique live unresolved bridge branch with the
alternate-$2$ entry state
\[
E=(m-1,1,u_{\mathrm{source}},2)
\]
on the canonical seed \(\CanSeed\).
The front-end one-corner reduction theorem
(Theorem~\ref{thm:conditional-front-end}) then reduces every unresolved
front-end instance to this normalized best-seed channel.

At that point the best-seed one-corner structural theorem package applies
directly to the branch started at \(E\): it proves candidate/actual agreement up
to first exit, identifies the exact trigger family \(H_{L1}\), pins the
universal exits to \(T_{\mathrm{reg}}\) and \(T_{\mathrm{exc}}\), and rules out
patched current classes before exit.
These are precisely the four assertions of the present theorem.
\end{proof}

\begin{remark}[provenance of the odometer spine]
In the attached theorem chain, Theorem~\ref{thm:odometer-spine} is the
manuscript-order form of the structural first-exit result proved in the one-pass
odd-$m$ globalization package.
The supporting normal-form and entry notes are exactly the one-corner reduction
notes tied to the canonical seed and the late local provenance check on the
surfaced $\sigma_{23}$/$R1$ branch.
\end{remark}

\begin{corollary}[front-end channel classification]
\label{cor:channel-classification}
For every odd modulus $m>2$, every unresolved front-end instance reduces to the
best-seed channel starting from the canonical seed
\[
\CanSeed=[2,2,1]/[1,4,4]
\]
and enters that channel at the alternate-$2$ state
\[
E=(m-1,1,u_{\mathrm{source}},2).
\]
Equivalently, the live unresolved front end is the one-corner channel
$R1\to H_{L1}$ with the entry described in Theorem~\ref{fr:T4}.
\end{corollary}

\begin{proof}
This is exactly the combination of the front-end reduction theorem
(Theorem~\ref{thm:conditional-front-end}), the canonical-entry theorem
(Theorem~\ref{fr:T4}), and the odometer identification in
Theorem~\ref{thm:odometer-spine}.
\end{proof}

\begin{remark}
This is the point at which the front-end proof changes language completely.
Before \(T4\), one argues in the grid language of rows and cells.
After \(T4\), the proof is an exact odometer proof in the coordinates
\((q,w,u,\Theta)\).
\end{remark}

\section{Selector surgery and graph-side Hamilton branches}
\label{sec:selector-surgery}

The present draft is not only about the front end.
The broader odd-$m$, $d=5$ paper also needs the graph-side part of the
odometer-view story.

\subsection{Standing theorem packets after the front end}

The current manuscript now records the remaining theorem-side packets in their
explicit theorem-order form rather than as anonymous imports.

\begin{definition}[theorem-side anchor quotient]
Let \(X_m\) denote the exact marked carry-jump-to-wrap object on the canonical
best-seed channel.
Define the theorem-side quotient
\[
Q_{\mathrm{th}}:=(B,\beta),
\]
where \(B\) is the grouped base and \(\beta\) is the canonical clock.
\end{definition}

\begin{definition}[the corrected selector families]
Let \(\SelCorr\) denote the explicit corrected selector whose rows are:
clean constant rules on \(\Sigma=0,1\) and on \(\Sigma\notin\{2,3\}\),
the explicit uniform slice-\(2\) correction \(q^{(2)}_S\),
and the explicit slice-\(3\) correction \(q^{(3)}_T\).

Let \(\SelStar\) denote the explicit selector surgery obtained by replacing the
slice-\(2\) rule \(q^{(2)}_S\) of \(\SelCorr\) by the explicit surgery table
\(q^{(2,\star)}_S\) and leaving the slice-\(3\) rule unchanged.
\end{definition}

\begin{remark}[definition status of the named graph-side objects]
The symbols \(Q_{\mathrm{th}}\), \(\SelCorr\), and \(\SelStar\) are used here
as theorem-order names tied to the corrected packet notes.
They are not yet expanded in this manuscript as full literal rule tables.
For formalization, the actual row formulas and quotient maps still have to be
imported from the corresponding packet notes.
\end{remark}

\begin{theorem}[theorem-side $M2/M3$ anchor package]
\label{thm:M23}
For every odd modulus $m\ge 5$, the exact marked carry-jump-to-wrap object
exists theorem-side with coordinates equivalent to
\[
(\rho,\beta,q_0,\sigma)
\qquad\text{and}\qquad
(B,\beta).
\]
The quotient \(Q_{\mathrm{th}}=(B,\beta)\) is exact deterministic for the
current event map, retains grouped base, and factors canonically to
\((\beta,\delta)\).
\end{theorem}

\begin{proof}
The corrected theorem-side package proves that the exact marked
carry-jump-to-wrap object exists and that the coordinates
\[
(B,\beta)
\qquad\text{and}\qquad
(\rho,\beta,q_0,\sigma)
\]
are equivalent parameterizations of the same exact marked object.
It also proves the corrected factor theorem: any quotient that is exact
deterministic for the current event map and retains grouped base factors
canonically through
\[
(B,\beta)\cong(\rho,\beta,q_0,\sigma)\twoheadrightarrow(\beta,\delta).
\]

For the intended theorem-side quotient we take
\[
Q_{\mathrm{th}}=(B,\beta).
\]
The quotient-verification step proves directly that this quotient is exact
deterministic for the current event map and retains grouped base.
Applying the corrected factor theorem to \(Q_{\mathrm{th}}\) yields the claimed
canonical factor chain, which is exactly the anchor package used later.
\end{proof}

\begin{theorem}[corrected-selector baseline theorem for $\SelCorr$]
\label{thm:M4-SelCorr}
For every odd modulus $m\ge 5$, the corrected selector family $\SelCorr$
defines five global torus permutations whose arc sets are pairwise disjoint and
cover the full directed arc set of the torus.
\end{theorem}

\begin{proof}
The direct corrected-selector theorem writes down the explicit slice-\(2\)
correction \(q^{(2)}_S\) and the explicit slice-\(3\) correction \(q^{(3)}_T\).
It proves incoming uniqueness on the only nonconstant target slices:
on target slice \(\Sigma=3\), the incoming color at generator \(j\) ignores the
single predecessor bit that can change at generator \(j\), so the whole
incoming row factors through one fixed subset \(M(y)-1\);
on target slice \(\Sigma=4\), only the transported bits
\(\mathbf 1_{1\in A_3}\) and \(\mathbf 1_{2\in A_3}\) matter, and they force the
incoming color pairs \(\{0,2\}\) and \(\{1,4\}\).
On every other slice the selector row is constant, so incoming uniqueness is
immediate there as well.

Hence each color map of \(\SelCorr\) has in-degree \(1\) and out-degree \(1\) at
every vertex, so each color map is a permutation of the torus.
At each vertex the selector row is a permutation of the five generators, so the
five selected arcs are exactly the five outgoing torus arcs.
Therefore the five arc sets are pairwise disjoint and cover the full directed
arc set.
\end{proof}

\begin{theorem}[color-$4$ Sel$^{\star}$ Hamilton theorem]
\label{thm:M5-color4}
For every odd modulus $m\ge 5$, the color-$4$ torus map coming from the surgery
family $\SelStar$ is Hamilton.
\end{theorem}

\begin{proof}
Let \(F_4^\star\) be the color-$4$ map induced by \(\SelStar\).
The selector-surgery route reduces Hamiltonicity of \(F_4^\star\) to the
third-return map on the final section.
The corrected-row stitching theorem proves that the explicit corrected-row model
\(U_m^\sharp\) on that final section is a single \(m^2\)-cycle for every odd
\(m\ge 11\).
The all-odd identification theorem then proves that the actual third return
\(U_m\) of \(F_4^\star\) agrees with \(U_m^\sharp\) for every odd \(m\ge 11\).
Therefore \(U_m\) is one cycle for every odd \(m\ge 11\), and the nested-return
lift criterion gives Hamiltonicity of \(F_4^\star\) on the full torus.

The remaining odd moduli \(m=5,7,9\) are covered by the direct exact checks
already recorded with the selector surgery and promoted in the all-odd
identification step.
Thus the color-$4$ torus map of \(\SelStar\) is Hamilton for every odd
\(m\ge 5\).
\end{proof}

\begin{theorem}[odd-$m$ D5 globalization theorem]
\label{thm:M6-globalization}
For every odd modulus $m\ge 5$, every actual lift of the exceptional cutoff row
$\delta=3m-3$ continues through
\[
3m-2\to 3m-1
\]
and lies in the regular continuing endpoint class there.
Consequently no mixed-status realized $\delta$ exists, fixed realized $\delta$
determines tail length and endpoint class, and raw global $(\beta,\delta)$ is
exact on the true accessible boundary union.
\end{theorem}

\begin{proof}
The one-pass odd-$m$ globalization package proves, in theorem order, the
structural first-exit theorem, the fixed-\(\delta\) tail-length reduction, the
chart/interface landing theorem, the regular closure theorem, the actual
exceptional interface landing theorem, and finally the odd-$m$ D5 globalization
theorem.
Its final theorem is exactly the present statement: every actual lift of the
exceptional cutoff row \(\delta=3m-3\) glues through
\[
3m-2\to 3m-1
\]
into the regular continuing endpoint class.
The listed consequences are exactly the consequences extracted there: no
mixed-status realized \(\delta\), fixed realized \(\delta\) determines tail
length and endpoint class, and raw global \((\beta,\delta)\) is exact on the
true accessible boundary union.
\end{proof}

\begin{remark}[packet status after 253]
At this point the remaining formal-target packets isolated by the
\texttt{253} closure map are all explicit theorem-order objects:
the post-entry odometer spine, the theorem-side anchor package \(M23\), the
corrected-selector baseline theorem \(M4\), the color-$4$ \(\SelStar\) branch
\(M5\), and the odd-$m$ globalization theorem \(M6\).
What remains later in the manuscript is the short downstream assembly together
with one graph-side compatibility question:
the explicit two-swap package of \(G1\) must be tied explicitly to the
color-relative hypothesis used by the abstract cyclic transport theorem \(G2\).
\end{remark}

\begin{remark}
The right graph-side question is no longer whether one can recover one common
selector family from scratch.
The correct question is whether the already closed color-$4$ branch can be
inserted into a colorwise selector package and then transported to the other
colors.
\end{remark}

\subsection{The graph-side theorem block}

\begin{definition}[colorwise selector package on the torus]
A family of local color rules
\[
g_0,\dots,g_4
\]
on the directed torus is called a \emph{colorwise selector package} when:
\begin{enumerate}[label=\textup{(\arabic*)},leftmargin=2.8em]
\item for every vertex, the five selected arcs are exactly the five outgoing
      torus arcs from that vertex;
\item for each color $i\in\{0,1,2,3,4\}$ and each vertex $x$, there is exactly
      one selected color-$i$ arc ending at $x$.
\end{enumerate}
\end{definition}

\begin{theorem}[G1: colorwise selector-compatibility theorem]
\label{fr:G1}
There exist five explicit local color rules
\[
g_0,\dots,g_4
\]
such that
\begin{enumerate}[label=\textup{(\arabic*)},leftmargin=2.8em]
\item outgoing exhaustiveness holds pointwise;
\item incoming uniqueness holds colorwise;
\item the induced global maps are torus permutations; and
\item \(g_4\) is exactly the already closed color-$4$ rule from the
      \(\SelStar\) branch.
\end{enumerate}
In fact one may take the explicit baseline package \(B\) from the graph-side
baseline theorem, define
\[
D_1=\{x:\sigma(x)=2,\ x_0=m-1\},
\]
and let \(D_2\) be the color-$4$ layer-$3$ defect union
\(L3^{(0,p=1)}\cup L3^{(1,p=0)}\) (equivalently
\(\sigma(x)=3\) together with the corresponding \((x_0,x_1+x_3)\) condition,
and \(D_2=\varnothing\) when \(m=3\)).
Then \(g_0=B_0\), \(g_3=B_3\), colors \(1\) and \(4\) are swapped on \(D_1\),
and colors \(2\) and \(4\) are swapped on \(D_2\).
\end{theorem}

\begin{proof}
This is the explicit two-swap splice of d5\_247.
The surfaced color-$4$ mixed branch equals \(B_1\) on \(D_1\), \(B_2\) on \(D_2\),
and \(B_4\) elsewhere.
The sets \(D_1\) and \(D_2\) are invariant under the translations
\(e_1-e_4\) and \(e_2-e_4\), respectively, so
\[
B_1(D_1)=B_4(D_1),
\qquad
B_2(D_2)=B_4(D_2).
\]
Hence the finite-defect splice criterion preserves the permutation property for
colors \(1,2,4\), while colors \(0,3\) remain baseline permutations.
Pointwise outgoing exhaustiveness is preserved because on \(D_1\) one merely
transposes the baseline directions of colors \(1\) and \(4\), and on \(D_2\) one
transposes the baseline directions of colors \(2\) and \(4\).
Thus the resulting package is colorwise permutation, pointwise outgoing-exhaustive,
and has the required exact color-$4$ branch.
\end{proof}

\begin{theorem}[G2 closes by cyclic transport from a colorwise package]
\label{thm:G2-closed}
Assume a colorwise package \(g_0,\dots,g_4\) on the full torus, and assume the
package is defined in color-relative
coordinates so that the color-$c$ rule is the cyclic shift of the color-$0$
rule.  Then every conjugacy-invariant property of the closed color-$4$ branch
transports to the remaining colors.  In particular, permutation, first-return,
grouped-return, nested-return, and Hamilton-type conclusions transfer from
color $4$ to colors \(0,1,2,3\).
\end{theorem}

\begin{proof}
Write \(\rho_c\) for cyclic coordinate shift by \(c\).  For a color-relative
local rule, the color-$c$ global map has the form
\[
 g_c = \rho_{-c}\circ g_0\circ \rho_c.
\]
Hence all five color maps are pairwise conjugate.  The listed graph-side
properties are invariant under conjugacy and under the corresponding transport
of the chosen section/return coordinates.
\end{proof}

\begin{corollary}[the explicit splice gives a colorwise selector package]
\label{cor:colorwise-package}
The explicit rules $g_0,\dots,g_4$ from Theorem~\ref{fr:G1} form a colorwise
selector package on the full torus.
\end{corollary}

\begin{proof}
The outgoing-exhaustiveness and colorwise incoming-uniqueness conclusions in
Theorem~\ref{fr:G1} are exactly the two axioms in the definition above.
\end{proof}

\begin{corollary}[all five colors are Hamilton, assuming transport compatibility]
\label{cor:all-five-hamilton}
Assume the explicit package of Theorem~\ref{fr:G1} is obtained from one
color-relative local rule family in the sense of
Theorem~\ref{thm:G2-closed}.
Then for every odd modulus $m\ge 5$, each global map induced by the package of
Theorem~\ref{fr:G1} is Hamilton.
\end{corollary}

\begin{proof}
Theorem~\ref{fr:G1} identifies $g_4$ with the already closed color-$4$ rule
coming from the $\SelStar$ branch, so Theorem~\ref{thm:M5-color4} makes $g_4$
Hamilton.  Theorem~\ref{thm:G2-closed} transports the same Hamilton-type
conclusion to colors $0,1,2,3$.
\end{proof}

\begin{remark}[graph-side status after the splice closure]
The old requirement of one common selector family is stronger than what the
final theorem actually needs.
What remained was exactly one compatibility theorem: splice the closed color-$4$
branch into a colorwise package while preserving pointwise outgoing
exhaustiveness and colorwise incoming uniqueness.
That selector-compatibility theorem is now closed by the explicit two-swap
package above.
What is still separate is the transport-compatibility step showing that this
explicit package also satisfies the color-relative hypothesis used by
Theorem~\ref{thm:G2-closed}.
\end{remark}

\section{Final upgrade and main theorem}

The endgame is now the corrected generic odd-$m$, $d=5$ upgrade step.
Compared with the older formulation, the corrected criterion does \emph{not}
ask for a single common selector family and does \emph{not} use the false
future-word separation target for the current event map.
Instead it asks for:
\begin{enumerate}[label=\textup{(\arabic*)},leftmargin=2.8em]
\item the best-seed channel classification on the front end;
\item the theorem-side anchor/globalization package;
\item a colorwise selector package on the full torus;
\item Hamiltonicity of each color map.
\end{enumerate}

\begin{theorem}[corrected generic odd-$m$, $d=5$ upgrade criterion]
\label{thm:generic-upgrade}
Let $m\ge 5$ be odd.
Assume:
\begin{enumerate}[label=\textup{(G\arabic*)},leftmargin=3.0em]
\item the front-end channel classification of
      Corollary~\ref{cor:channel-classification};
\item the theorem-side anchor package of Theorem~\ref{thm:M23} together with
      the globalization theorem~\ref{thm:M6-globalization};
\item a colorwise selector package on the full torus;
\item Hamiltonicity of each color map in that package.
\end{enumerate}
Then the directed torus $D_5(m)$ admits a decomposition into five arc-disjoint
Hamilton cycles.
\end{theorem}

\begin{proof}
By (G3), the five color maps determine five selected arc sets whose outgoing
arcs at each vertex are exactly the five torus outgoing arcs and whose incoming
arcs are unique colorwise.  Hence the five global maps are torus permutations
with pairwise disjoint arc sets covering the full directed arc set.
By (G4), each of those permutations is a Hamilton cycle.
Therefore the five arc sets form a Hamilton decomposition of $D_5(m)$.
The role of (G1) and (G2) is not additional combinatorics at this last step;
it is to identify the selector package and its Hamiltonicity with the actual
odd-$m$, $d=5$ branch selected by the front-end and theorem-side packets.
\end{proof}

\begin{proof}[Proof of Theorem~\ref{thm:conditional-full}]
Fix an odd modulus $m\ge 5$.
The front-end theorem block closes the seed search and leaves a unique live
channel; Corollary~\ref{cor:channel-classification} records this as the
best-seed one-corner channel.
The post-entry dynamics on that channel are then governed by the explicit
odometer spine of Theorem~\ref{thm:odometer-spine}.
On the theorem side, Theorem~\ref{thm:M23} supplies the exact anchor package
and Theorem~\ref{thm:M6-globalization} supplies the odd-$m$ globalization
conclusion.
On the graph side, Theorem~\ref{fr:G1} and
Corollary~\ref{cor:colorwise-package} give a full colorwise selector package,
while Corollary~\ref{cor:all-five-hamilton} gives Hamiltonicity of all five
colors under the standing transport-compatibility hypothesis in the statement
of the theorem.
Therefore all hypotheses of Theorem~\ref{thm:generic-upgrade} are satisfied,
and the directed torus $D_5(m)$ admits a decomposition into five arc-disjoint
directed Hamilton cycles.
\end{proof}

\begin{remark}
The final theorem is now reduced to one explicit remaining graph-side
compatibility input plus the already packaged theorem blocks.
What remains for later work is therefore not another broad theorem search inside
the odd-$m$, $d=5$ chain, but either:
\begin{enumerate}[label=\textup{(\arabic*)},leftmargin=2.8em]
\item make the transport compatibility of the explicit \(G1\) package fully
      explicit, or
\item keep that step as a standing named input in the final theorem.
\end{enumerate}
\end{remark}

\section{Dependency ledger for formalization}

For later Leanization or line-level polishing, the main dependency chain used in
Theorem~\ref{thm:conditional-full} is the following.
\begin{enumerate}[label=\textup{(D\arabic*)},leftmargin=3.0em]
\item finite seed certificate and bridge normalization:
      Proposition~\ref{prop:finite-chart} and
      Proposition~\ref{prop:bridge-normalization};
\item the actual-needed tiny-repair core:
      Theorem~\ref{thm:T0-needed};
\item the row-theorem block:
      Theorems~\ref{thm:T2-closed}, \ref{thm:T3-closed},
      \ref{thm:sigma23-closed}, and \ref{thm:sigma32-closed}, summarized in
      Theorem~\ref{thm:conditional-front-end};
\item canonical entry and one-corner odometer identification:
      Theorem~\ref{fr:T4}, Theorem~\ref{thm:odometer-spine}, and
      Corollary~\ref{cor:channel-classification};
\item theorem-side anchor and globalization packets:
      Theorems~\ref{thm:M23} and \ref{thm:M6-globalization};
\item graph-side baseline and closed color-$4$ branch:
      Theorems~\ref{thm:M4-SelCorr} and \ref{thm:M5-color4};
\item selector splice and color transport:
      Theorem~\ref{fr:G1}, Theorem~\ref{thm:G2-closed},
      the realization of the \(G2\) hypothesis for the explicit splice package,
      Corollaries~\ref{cor:colorwise-package} and
      \ref{cor:all-five-hamilton};
\item corrected final upgrade:
      Theorem~\ref{thm:generic-upgrade};
\item main odd-$m$, $d=5$ theorem:
      Theorem~\ref{thm:conditional-full}.
\end{enumerate}

\begin{remark}[status after closure]
At the level of packet closure, no live external packet remains among the
objects isolated by the \(253\) closure map.
The front-end slice $T0$--$T4$ is closed in the present bundle, the packet notes
\(M23\), \(M4\), \(M5\), \(M6\) are now written as explicit theorem-order
objects, and the remaining standing inputs are named precisely.
The main remaining tightening point for a fully self-contained final theorem is
the graph-side transport compatibility linking the explicit \(G1\) package to
the color-relative hypothesis of \(G2\).
\end{remark}

\section{Heuristic picture: bridge decentering}

\begin{heuristic}[bridge decentering]
The front-end chart suggests the following picture.
Leading and centered bridge positions are unstable.
The front end tries to decenter toward bridge-compatible terminal positions,
except that the repeated-end terminal corner \(\sig33\) is itself obstructed.
This heuristic has now been realized by local slice theorems rather than by a
single global monotonicity theorem:
\begin{itemize}[leftmargin=2.4em]
\item first-column companion or centered-strip escape for \(T1\);
\item centered-core incompatibility for \(T2\);
\item repeated-end obstruction for \(T3\).
\end{itemize}
\end{heuristic}

\begin{remark}
This heuristic is useful for explanation and proof search, but it is not itself
a theorem.
The actual manuscript should continue to use local theorem objects rather than a
single global potential story, unless a genuinely odd-uniform monotonicity
theorem is later found.
\end{remark}

\section{Editorial note}

The mathematical chain recorded in this manuscript is now stable.
Future passes should therefore focus on:
\begin{enumerate}[label=\textup{(\arabic*)},leftmargin=2.8em]
\item line-level exposition cleanup;
\item theorem naming and local notation polishing;
\item optional further internalization of older upstream packets;
\item synchronization with any Lean formalization ledger.
\end{enumerate}

\begin{remark}
The important point is that these are editorial or packaging tasks.
They are no longer theorem-search tasks inside the odd-$m$, $d=5$ proof chain
used for Theorem~\ref{thm:conditional-full}.
\end{remark}

\end{document}
